\chapter{字符串哈希}

% \begin{overview}
% \begin{itemize}
%     \item \texttt{vector}、栈、队列、链表、堆支持的操作，及其实现
%     \item 常见的可以使用这些数据结构解决的问题
% \end{itemize}
% \end{overview}

\section{哈希函数及其碰撞概率}

字符串哈希是比较字符串最简单的方法。通常来说，我们使用如下的方式将字符串映射到 $\mathbb F_p$ 中的一个元素：

\begin{descbox}{算法}[字符串哈希]
    假设字符集 $\Sigma=\{1,\ldots,c\}$。对于素数 $p>|\Sigma|$ 和整数参数 $0\le x<p$，字符串哈希函数定义为
    $$
    f_x(S)=\left(\sum_{i=1}^{|S|} S[i]\cdot x^{|S|-i}\right) \bmod p
    $$.

    在使用时，随机选择 $0\le x<p$，然后将全局的哈希函数固定为 $f=f_x$。
\end{descbox}
\begin{remark}
    字符集不包含 $0$ 是重要的，否则无法区分不同长度的字符串。
    
    在 OI 中，$x$ 常常会用一个\textbf{出题人猜不到的}固定常数替代随机数。

    哈希函数也可以反过来定义为
    $$
    f_x(S)=\left(\sum_{i=0}^{|S|-1} S[i+1]\cdot x^{i}\right) \bmod p.
    $$
\end{remark}

如果 $S,T$ 是相同的字符串，那么它们映射到同一个值；如果是不同的字符串，则它们大概率映射到不同的值。下面我们分析这个哈希函数的碰撞概率。

\begin{theorem}
    对任意两个不同的字符串 $S\neq T$，有 $\Pr[f(S)=f(T)]\le \frac{\max(|S|,|T|)-1}{p}$。
\end{theorem}
\begin{proof}
    记 $l=\max(|S|,|T|)$。注意到 $f(S)$ 和 $f(T)$ 都是不超过 $l-1$ 次的关于 $x$ 的多项式，$f(S)=f(T)$ 即是说这两个多项式的差为 $0$。由于 $S\ne T$，$f(S)-f(T)$ 是一个非零多项式，而 $\mathbb F_p$ 中的 $d$ 次非零多项式至多有 $d$ 个根，故 $x$ 落在这个多项式的根上的概率不超过 $\frac{l-1}{p}$，即 $\Pr[f(S)=f(T)]\le \frac{\max(|S|,|T|)-1}{p}$。
\end{proof}
\begin{corollary}
    对于 $n$ 个长度不超过 $l$ 的字符串，哈希函数 $f$ 发生碰撞的概率不超过 $\frac{n(n-1)(l-1)}{2p}$。
\end{corollary}


在 OI 中 $p$ 通常在 $10^9$ 级别，而 $n,L$ 常常达到 $10^5$ 级别，此时的碰撞概率已经非常大。为了避免这个问题，除了使用 $10^{18}$ 级别的模数以外，更常见的做法是\textbf{双哈希}，即使用两个独立的哈希函数 $f_x,f_y$。此时，两个串的碰撞概率降低至 $(l/p)^2$。

\section{子串哈希}

我们的哈希函数可以方便地提取 $S$ 的某个子串 $S[l,r]$ 的哈希值：
$$
\begin{aligned}
f(S[l,r])&=\left(\sum_{i=l}^r S[i]\cdot x^{r-i}\right) \bmod p\\
&=\left(\sum_{i=1}^{r} S[i]\cdot x^{r-i}-x^{r-l+1}\sum_{i=1}^{l-1} S[i]\cdot x^{l-1-i}\right) \bmod p\\
\end{aligned}
$$
于是我们只要预处理所有的前缀哈希值以及 $x$ 的幂，就可以在 $O(1)$ 时间内计算任意子串的哈希值。

对于字符串有修改的情况，可以使用数据结构维护。

\section{应用}

以下是字符串哈希的常见应用。假设哈希值均已预处理完成。

\begin{itemize}
    \item 比较两个字符串是否相同
    \item $O(\log n)$ 查询两个字符串的最长公共前缀/后缀
    \item $O(\log n)$ 比较两个字符串的字典序
    \item $O(1)$ 判断字符串是否回文
    \item $O(n\log n)$ 计算一个字符串的最长回文子串
    \item $O(n\log n)$ 查询两个串的最长公共子串，$n$ 是两串长度的较小值
\end{itemize}

直接比较两个字符串的哈希值即可判断它们是否相等。也可以进一步用于比较两个字符串的字典序：二分出最长的相同前缀，然后比较下一位即可。

% \begin{summary}
% \end{summary}
